\documentclass[11pt]{article}

\usepackage[T1]{fontenc}
\usepackage{amsmath,amssymb,amsthm}
\usepackage[margin=1in]{geometry}
\usepackage[colorlinks=true,linkcolor=blue,urlcolor=blue]{hyperref}

\newtheorem{theorem}{Theorem}
\newtheorem{proposition}[theorem]{Proposition}
\newtheorem{corollary}[theorem]{Corollary}
\newtheorem{conjecture}[theorem]{Conjecture}
\theoremstyle{remark}
\newtheorem{remark}[theorem]{Remark}

\newcommand{\R}{\mathbb R}
\newcommand{\T}{\mathbb T}
\newcommand{\Pp}{\mathbb P}
\newcommand{\norm}[1]{\left\|#1\right\|}
\newcommand{\ip}[2]{\left\langle #1,#2\right\rangle}
\newcommand{\Var}{\operatorname{Var}}
\setlength{\emergencystretch}{2em}

\title{Radial Triad Spectrum and Dynamic Shell Descent\\
at the q4 Nonlinear Entrance}
\author{John Lawson and Codex}
\date{25 July 2026}

\begin{document}
\maketitle

\begin{abstract}
This report continues the conditional energy-efficient exact q4 branch.
The preceding round expressed half of the entrance energy as a diagonal
Gaussian nonlinear flux.  Here we identify its exact signed radial
spectrum and couple it to the positive radial kinetic-energy spectrum.
If \(\mathcal E(r,s)=\frac12\norm{e^{\nu r\Delta}v(s)}_2^2\), then
\[
 \Pi_r(s)=(\partial_s-\partial_r)\mathcal E(r,s).
\]
Thus every proposed flux must reconstruct to a function completely
monotone in heat age.  This excludes the preceding Laguerre scalar
survivor: its reconstructed energy has a sign-changing generalized
Laguerre inverse density.

The positivity correction does not close q4.  We construct a moving
positive spectral shell whose exact characteristic flux has profile
\(xe^{-x}\), satisfies the energy, diagonal, and recycling identities,
matches the q4 power--log \(M^2\) and enstrophy tails, saturates the
divergent \(MY^2\) variation budget, and places asymptotically all work
in shared early history.  An exact Pell family of divergence-free torus
triads transfers energy between adjacent radial shells with positive
heat flux for every heat age; after parabolic rescaling its profile
converges to \((\sqrt3/8)xe^{-x}\).

The remaining gate is consequently temporal rather than instantaneous:
the positive shell motion and adjacent-shell quadratic transfer must be
identified along one Navier--Stokes trajectory.  The shell is not a
velocity field and the Pell family is not one evolution.  All results
are conditional repository deductions pending external review.  No Clay
alternative is proved.
\end{abstract}

\section{Conditional entrance and radial measures}

Let \(v,c,U=v+c\) be the conditional reversed full-defect entrance from
the preceding rounds, and set
\begin{equation}
\label{eq:F}
 F(s):=\Pp((U\cdot\nabla)v)(s).
\end{equation}
It obeys
\begin{equation}
\label{eq:evolution}
 \partial_s v-\nu\Delta v+F=0,
 \qquad \nabla\cdot U=\nabla\cdot v=0,
\end{equation}
\begin{equation}
\label{eq:energy}
 \norm{v(s)}_2^2+
 2\nu\int_0^s\norm{\nabla v(a)}_2^2\,da=d>0,
\end{equation}
and, for each fixed \(r>0\),
\begin{equation}
\label{eq:heat-boundary}
 e^{\nu r\Delta}v(s)\longrightarrow0
 \quad\hbox{strongly in \(L^2\) as \(s\downarrow0\)}.
\end{equation}
The diagonal flux is
\begin{equation}
\label{eq:Pi}
 \Pi_r(s):=-\ip{F(s)}{e^{2\nu r\Delta}v(s)}.
\end{equation}

Use the unitary Fourier transform
\[
 \widehat f(\xi)
 =(2\pi)^{-3/2}\int_{\R^3}e^{-ix\cdot\xi}f(x)\,dx
\]
and push forward by
\begin{equation}
\label{eq:lambda}
 \lambda=2\nu|\xi|^2.
\end{equation}
Define a positive measure \(\varepsilon_s\) and a signed measure
\(\mu_s\) by
\begin{align}
\label{eq:epsilon}
 \varepsilon_s(B)
 &:=
 \frac12\int_{\{2\nu|\xi|^2\in B\}}
 |\widehat v(s,\xi)|^2\,d\xi,\\
\label{eq:mu}
 \mu_s(B)
 &:=
 -\operatorname{Re}
 \int_{\{2\nu|\xi|^2\in B\}}
 \widehat F(s,\xi)\cdot
 \overline{\widehat v(s,\xi)}\,d\xi.
\end{align}

\begin{theorem}[Exact radial spectra]
\label{thm:radial}
For
\begin{equation}
\label{eq:heat-energy}
 \mathcal E(r,s):=
 \frac12\norm{e^{\nu r\Delta}v(s)}_2^2,
\end{equation}
one has
\begin{align}
\label{eq:E-Laplace}
 \mathcal E(r,s)
 &=\int_0^\infty e^{-r\lambda}\,d\varepsilon_s(\lambda),\\
\label{eq:Pi-Laplace}
 \Pi_r(s)
 &=\int_0^\infty e^{-r\lambda}\,d\mu_s(\lambda).
\end{align}
If the Fourier pairings have densities, then
\begin{equation}
\label{eq:radial-density}
\begin{aligned}
 \frac{d\mu_s}{d\lambda}(\lambda)
 =
 -\frac1{4\nu}\sqrt{\frac{\lambda}{2\nu}}\,
 \operatorname{Re}\int_{\mathbb S^2}
 &\widehat F
 \left(s,\sqrt{\frac{\lambda}{2\nu}}\omega\right)\\
 &\cdot
 \overline{\widehat v
 \left(s,\sqrt{\frac{\lambda}{2\nu}}\omega\right)}
 \,d\omega .
\end{aligned}
\end{equation}
Moreover,
\begin{equation}
\label{eq:mass-zero}
 \mu_s([0,\infty))=0.
\end{equation}
\end{theorem}

\begin{proof}
Plancherel, self-adjointness, and the heat semigroup law give
\[
 \mathcal E(r,s)
 =
 \frac12\int_{\R^3}
 e^{-2\nu r|\xi|^2}|\widehat v(s,\xi)|^2\,d\xi,
\]
which is \eqref{eq:E-Laplace}.  The same calculation applied to
\eqref{eq:Pi} gives \eqref{eq:Pi-Laplace}.  In polar coordinates,
\[
 |\xi|^2\,d|\xi|
 =
 \frac1{4\nu}\sqrt{\frac{\lambda}{2\nu}}\,d\lambda,
\]
proving \eqref{eq:radial-density}.  Finally,
\[
 \mu_s([0,\infty))
 =-\ip{F(s)}{v(s)}
 =-\int_{\R^3}(U\cdot\nabla)v\cdot v\,dx=0
\]
by incompressibility.
\end{proof}

\begin{proposition}[Quadratic triad density]
\label{prop:triad-density}
Let \(c_0=(2\pi)^{-3/2}\).  Before radial pushforward, the signed density
in \eqref{eq:mu} is
\begin{equation}
\label{eq:triad-density}
\begin{aligned}
 -\operatorname{Re}
 \bigl(\widehat F(s,\xi)\cdot\overline{\widehat v(s,\xi)}\bigr)
 =
 c_0\operatorname{Im}\int_{\R^3}
 &\bigl((\xi-\eta)\cdot\widehat U(s,\eta)\bigr)\\
 &\times
 \bigl(\widehat v(s,\xi-\eta)\cdot
 \overline{\widehat v(s,\xi)}\bigr)\,d\eta .
\end{aligned}
\end{equation}
\end{proposition}

\begin{proof}
Fourier transformation gives
\[
 \widehat F_i(\xi)
 =
 ic_0P_{i\ell}(\xi)
 \int_{\R^3}
 (\xi-\eta)_j\widehat U_j(\eta)
 \widehat v_\ell(\xi-\eta)\,d\eta.
\]
The Leray projection disappears when paired with
\(\overline{\widehat v(\xi)}\), since the latter is perpendicular to
\(\xi\).  The identity
\(-\operatorname{Re}(iz)=\operatorname{Im}z\) completes the proof.
\end{proof}

\begin{remark}
Pressure supplies no radial sign.  The split \(U=v+c\) splits \(\mu_s\)
into self and drift measures, each of total mass zero.  Strong \(L^2\)
convergence of \(c\) alone gives no clock-relative control of higher
radial moments.
\end{remark}

\section{The spectral characteristic law}

\begin{theorem}[Positive heat energy along characteristics]
\label{thm:characteristic}
For every positive smooth time,
\begin{equation}
\label{eq:characteristic}
 \boxed{
 \Pi_r(s)=(\partial_s-\partial_r)\mathcal E(r,s).
 }
\end{equation}
Equivalently, in radial spectral distributions,
\begin{equation}
\label{eq:measure-evolution}
 \boxed{
 \mu_s=\partial_s\varepsilon_s+\lambda\varepsilon_s.
 }
\end{equation}
For every integer \(k\ge0\),
\begin{equation}
\label{eq:CM}
 \boxed{
 (-1)^k\partial_r^k\mathcal E(r,s)
 =
 \int_0^\infty
 \lambda^ke^{-r\lambda}\,d\varepsilon_s(\lambda)
 \ge0.
 }
\end{equation}
Finally, the boundary condition \eqref{eq:heat-boundary} gives
\begin{equation}
\label{eq:reconstruction}
 \boxed{
 \mathcal E(r,s)
 =
 \int_0^s\Pi_{r+s-a}(a)\,da.
 }
\end{equation}
\end{theorem}

\begin{proof}
Differentiating \eqref{eq:heat-energy} in \(r\) gives
\begin{equation}
\label{eq:E-r}
 \partial_r\mathcal E(r,s)
 =-\nu\norm{\nabla e^{\nu r\Delta}v(s)}_2^2.
\end{equation}
Using \eqref{eq:evolution},
\begin{equation}
\label{eq:E-s}
\begin{aligned}
 \partial_s\mathcal E(r,s)
 &=
 \ip{\nu\Delta v-F}{e^{2\nu r\Delta}v}\\
 &=
 -\nu\norm{\nabla e^{\nu r\Delta}v}_2^2+\Pi_r(s).
\end{aligned}
\end{equation}
Subtracting \eqref{eq:E-r} from \eqref{eq:E-s} proves
\eqref{eq:characteristic}; inverse Laplace transformation proves
\eqref{eq:measure-evolution}.  Repeated differentiation of
\eqref{eq:E-Laplace} proves \eqref{eq:CM}.

Follow the characteristic \((r+s-a,a)\), \(0\le a\le s\).  Equation
\eqref{eq:characteristic} gives
\[
 \frac d{da}\mathcal E(r+s-a,a)=\Pi_{r+s-a}(a).
\]
Its lower endpoint vanishes by \eqref{eq:heat-boundary}.  Integration
proves \eqref{eq:reconstruction}.
\end{proof}

\begin{corollary}[Flux admissibility audit]
\label{cor:audit}
Every proposed diagonal flux must reconstruct through
\eqref{eq:reconstruction} to a completely monotone heat-age function.
A signed Laplace representation of the flux and the cancellation
\(\Pi_0=0\) are not sufficient.
\end{corollary}

\section{The Laguerre survivor is not positive energy}

Freeze the limiting energy at \(d\) and define
\begin{equation}
\label{eq:old-Pi}
 \widetilde\Pi_n(r,s)
 :=
 \frac{(n+1)d}{2}\frac{r^n}{(r+s)^{n+1}},
 \qquad n\ge1.
\end{equation}

\begin{proposition}[Sign-changing reconstructed spectrum]
\label{prop:laguerre-fails}
The characteristic reconstruction of \eqref{eq:old-Pi} is
\begin{equation}
\label{eq:old-E}
 \widetilde{\mathcal E}_n(r,s)
 =
 \frac d2
 \left[
 1-\left(\frac r{r+s}\right)^{n+1}
 \right].
\end{equation}
Its inverse Laplace density is
\begin{equation}
\label{eq:old-density}
 \frac d2\,s e^{-s\lambda}L_n^{(1)}(s\lambda),
\end{equation}
which changes sign.  Hence \eqref{eq:old-E} is not completely monotone.
\end{proposition}

\begin{proof}
Along the characteristic, the two arguments of
\(\widetilde\Pi_n\) sum to \(r+s\), so
\[
\begin{aligned}
 \int_0^s\widetilde\Pi_n(r+s-a,a)\,da
 &=
 \frac{(n+1)d}{2(r+s)^{n+1}}
 \int_0^s(r+s-a)^n\,da\\
 &=\widetilde{\mathcal E}_n(r,s).
\end{aligned}
\]
Expanding the bracket yields
\[
 \widetilde{\mathcal E}_n(r,s)
 =
 \frac d2\sum_{k=1}^{n+1}
 (-1)^{k+1}\binom{n+1}{k}
 \frac{s^k}{(r+s)^k}.
\]
Since
\[
 (r+s)^{-k}
 =
 \frac1{(k-1)!}\int_0^\infty
 \lambda^{k-1}e^{-(r+s)\lambda}\,d\lambda,
\]
the inverse density is
\[
 \frac d2\,s e^{-s\lambda}
 \sum_{j=0}^n(-1)^j
 \binom{n+1}{j+1}\frac{(s\lambda)^j}{j!},
\]
which is \eqref{eq:old-density}.  The generalized Laguerre polynomial
\(L_n^{(1)}\) has \(n\) positive simple roots and changes sign.
Bernstein's theorem therefore rules out complete monotonicity.
\end{proof}

\begin{remark}
The full preceding survivor multiplies the same core by a slowly varying
factor \(E_\sharp(r+s)\).  Under \(r=sx\), that factor converges with
rescaled derivatives to \(d\), with power--log error \(O(s^{2/11})\).
The strict derivative-sign failure persists at small \(s\).  The old
kernel remains a countermodel to the scalar estimates of its round, but
not to positive kinetic spectral energy.
\end{remark}

\section{A positive q4 moving-shell survivor}

Put
\begin{equation}
\label{eq:L}
 \ell(s):=\log(e/s),
 \qquad
 L(s):=s^{-9/11}\ell(s)^{2/11},
\end{equation}
\begin{equation}
\label{eq:A}
 A(s):=\frac d2
 \exp\left(-\int_0^sL(a)\,da\right),
\end{equation}
and
\begin{equation}
\label{eq:shell-E}
 \mathcal E_\circ(r,s):=A(s)e^{-rL(s)}.
\end{equation}
Define
\begin{equation}
\label{eq:a}
 a(s):=-\frac{sL'(s)}{L(s)}
 =\frac9{11}+\frac{2}{11\ell(s)}.
\end{equation}

\begin{theorem}[Positive-energy q4 reuse ledger]
\label{thm:shell}
The shell energy \eqref{eq:shell-E} is completely monotone, vanishes at
every fixed positive heat age as \(s\downarrow0\), and has characteristic
flux
\begin{equation}
\label{eq:shell-Pi}
 \boxed{
 \Pi^\circ_r(s)
 =
 -A(s)rL'(s)e^{-rL(s)}
 =
 \frac{A(s)a(s)}s
 \bigl(rL(s)\bigr)e^{-rL(s)}>0.
 }
\end{equation}
It satisfies
\begin{equation}
\label{eq:shell-energy}
 A(s)+\int_0^sA(a)L(a)\,da=\frac d2.
\end{equation}
With
\begin{equation}
\label{eq:shell-MY}
 Y_\circ(s)^2:=\frac{A(s)L(s)}{\nu},
 \qquad
 M_\circ(s):=\frac1{sL(s)},
\end{equation}
one has
\begin{align}
\label{eq:M-tail}
 \int_0^hM_\circ^2
 &\asymp h^{7/11}\ell(h)^{-4/11},\\
\label{eq:Y-tail}
 \int_0^hY_\circ^2
 &\asymp h^{2/11}\ell(h)^{2/11},\\
\label{eq:MY2}
 M_\circ(s)Y_\circ(s)^2
 &=\frac{A(s)}{\nu s}.
\end{align}
Furthermore,
\begin{align}
\label{eq:shell-bound}
 |\Pi^\circ_r(s)|
 &\lesssim r^{-1/2}M_\circ(s)Y_\circ(s),\\
\label{eq:shell-var}
 \Var_{r>0}\Pi^\circ_r(s)
 &=\frac{2A(s)a(s)}{e\,s}
 \asymp M_\circ(s)Y_\circ(s)^2.
\end{align}
It obeys the exact diagonal and recycling identities, and for every fixed
\(0<q<1\),
\begin{equation}
\label{eq:early}
 \boxed{
 \frac{\displaystyle
 \int_0^{q\tau}\Pi^\circ_{\tau-s}(s)\,ds}
 {\displaystyle
 \int_0^\tau\Pi^\circ_{\tau-s}(s)\,ds}
 \longrightarrow1
 \quad(\tau\downarrow0).
 }
\end{equation}
\end{theorem}

\begin{proof}
Equation \eqref{eq:shell-E} is the Laplace transform of the positive
measure \(A(s)\delta_{L(s)}\), so it is completely monotone.  Since
\(L(s)\to\infty\), its fixed positive heat-age limit is zero.  Direct
differentiation gives
\[
 A'(s)=-L(s)A(s)
\]
and \eqref{eq:shell-Pi}; integration gives
\eqref{eq:shell-energy}.  The characteristic identity then supplies
the diagonal and recycling laws.

The power--log integral
\[
 \int_0^sL(a)\,da
 \asymp s^{2/11}\ell(s)^{2/11}
\]
shows \(A(s)\to d/2\).  Equations \eqref{eq:L} and
\eqref{eq:shell-MY} give
\[
 M_\circ=s^{-2/11}\ell^{-2/11},
 \qquad
 Y_\circ^2\asymp s^{-9/11}\ell^{2/11},
\]
and one-sided power--log integration proves
\eqref{eq:M-tail}--\eqref{eq:MY2}.

Put \(x=rL(s)\).  Then
\[
 \Pi^\circ_r(s)=\frac{A(s)a(s)}sxe^{-x}
\]
and
\[
 r^{-1/2}M_\circ Y_\circ
 =
 \sqrt{\frac{A(s)}{\nu}}\frac{x^{-1/2}}s.
\]
Boundedness of \(x^{3/2}e^{-x}\) proves
\eqref{eq:shell-bound}.  The profile \(xe^{-x}\) rises from zero to
\(e^{-1}\) and returns to zero, proving \eqref{eq:shell-var}.

The denominator in \eqref{eq:early} equals \(A(\tau)\to d/2\).
On \(q\tau\le s\le\tau\), use \(xe^{-x}\le x\), \(s\asymp\tau\), and
monotonicity of \(L\):
\[
 \int_{q\tau}^{\tau}\Pi^\circ_{\tau-s}(s)\,ds
 \lesssim_q\tau L(q\tau)
 \lesssim_q
 \tau^{2/11}\ell(\tau)^{2/11}\longrightarrow0.
\]
This proves \eqref{eq:early}.  The dominant source time satisfies
\(\tau L(s_\tau)\asymp1\), hence \(s_\tau/\tau\to0\).
\end{proof}

\begin{remark}
The moving delta shell is a legitimate ideal positive energy spectrum.
Its time derivative contains a derivative of a delta, so its transfer
measure is distributional in \(\lambda\).  Smooth shell mollifications
retain the asymptotic conclusions.  No quadratic velocity or NSE
evolution is asserted.
\end{remark}

\section{Exact adjacent-shell Pell triads}

Let \(n,m\) be positive integers solving
\begin{equation}
\label{eq:Pell}
 n^2-3m^2=1.
\end{equation}
On \(\T^3\), define
\begin{equation}
\label{eq:W}
\begin{aligned}
 W_{n,m}(x,y,z)
 :=
 \bigl(
 &-m\sin(nx+my),\,
 n\sin(nx+my),\\
 &\sin(2my)+\sin(nx-my)
 \bigr).
\end{aligned}
\end{equation}

\begin{theorem}[Positive adjacent-shell quadratic heat flux]
\label{thm:Pell}
The field \(W_{n,m}\) is real, smooth, mean zero, and divergence-free.
For every radial Fourier multiplier \(M\), whose shell values are
\(\eta_j\),
\begin{equation}
\label{eq:Pell-flux}
 \boxed{
 \frac1{|\T^3|}\int_{\T^3}
 W_{n,m}\otimes W_{n,m}:\nabla(MW_{n,m})\,dx
 =
 \frac{nm}{2}
 \bigl(\eta_{4m^2}-\eta_{4m^2+1}\bigr).
 }
\end{equation}
For \(M=e^{2\nu r\Delta}\),
\begin{equation}
\label{eq:Pell-heat}
 \boxed{
 Q_{n,m}(r)
 =
 \frac{nm}{2}e^{-8\nu m^2r}
 \bigl(1-e^{-2\nu r}\bigr)>0
 \quad(r>0).
 }
\end{equation}
Along the Pell sequence \(n/m\to\sqrt3\),
\begin{equation}
\label{eq:Pell-limit}
 \boxed{
 Q_{n,m}\left(\frac{x}{8\nu m^2}\right)
 \longrightarrow\frac{\sqrt3}{8}xe^{-x}
 }
\end{equation}
locally uniformly for \(x\ge0\).
\end{theorem}

\begin{proof}
The Fourier waves are
\[
 \pm(0,-2m,0),\qquad
 \pm(n,m,0),\qquad
 \pm(-n,m,0).
\]
Their positive-wave coefficients are \(i/2\) times
\[
 (0,0,1),\qquad(m,-n,0),\qquad(0,0,1),
\]
respectively.  Each coefficient is perpendicular to its wave.  The
three positive representatives sum to zero, and their squared radii are
\[
 4m^2,\qquad n^2+m^2=4m^2+1,
\]
using \eqref{eq:Pell}.

Exact enumeration in
\[
 \sum_{k+p+q=0}
 \widehat W_\alpha(k)\widehat W_\beta(p)
 (iq_\beta)\eta_{|q|^2}\widehat W_\alpha(q)
\]
gives the two shell coefficients \(nm/2\) and \(-nm/2\).  This proves
\eqref{eq:Pell-flux}; the repository's exact-rational certificate
independently repeats the enumeration.  Substitution of
\(\eta_j=e^{-2\nu rj}\) proves \eqref{eq:Pell-heat}.

For \(r=x/(8\nu m^2)\),
\[
 Q_{n,m}(r)
 =
 \frac{nm}{2}e^{-x}
 \left(1-e^{-x/(4m^2)}\right).
\]
Since \(n/m\to\sqrt3\), the last display converges locally uniformly to
\((\sqrt3/8)xe^{-x}\).
\end{proof}

\begin{remark}
The signed nonlinear spectrum has two adjacent atoms, zero total mass,
and a nonzero first radial moment.  Thus three-dimensional
incompressibility gives neither a first-moment cancellation nor a
heat-age sign change.  Standard solenoidal cutoff localizes each fixed
Pell field to \(\R^3\) on bounded heat-age ranges, but does not make the
sequence one trajectory.
\end{remark}

\section{Round disposition}

\subsection*{Robust findings, subject to external review}

\begin{enumerate}
 \item The actual Gaussian flux has the signed radial spectrum
 \eqref{eq:mu} and explicit quadratic density
 \eqref{eq:triad-density}.
 \item It is the characteristic derivative of a positive completely
 monotone heat-energy spectrum.
 \item The previous Laguerre survivor fails this positive-energy audit.
 \item The moving-shell model restores the q4 power--log tails, energy
 and recycling identities, Gaussian bounds, divergent variation budget,
 and asymptotically complete early reuse.
 \item The Pell family gives exact adjacent-shell incompressible
 self-interactions with positive flux at every heat age.
 \item Its rescaled quadratic profile converges to the shell survivor's
 \(xe^{-x}\) shape.
\end{enumerate}

\subsection*{Closed shortcuts}

None of complete monotonicity alone, total nonlinear spectral mass,
a universal first-moment cancellation, instantaneous self-interaction
sign, pressure projection, incompressibility, or adjacent-shell
separation closes q4.  The original Laguerre kernel is closed, but the
positive moving shell replaces it.

\subsection*{Open obligations}

\begin{enumerate}
 \item Derive a shell-speed law for a fixed quantile of
 \(\varepsilon_s\) from \eqref{eq:triad-density}.
 \item Prove that repeated positive adjacent-shell transfer carries a
 finite same-trajectory genealogy cost that the moving-shell survivor
 violates.
 \item Quantify smooth shell thickness instead of the delta-shell
 idealisation.
 \item Exclude coherent steering by the strong-trace drift.
 \item Couple the radial shell descent to the Hausdorff-one-null boundary
 energy measure and spatial capacity.
 \item Submit the conditional chain to independent external review.
 \item Treat slower clocks, divergent normalised energy, R3B, and the
 other Clay alternatives separately.
\end{enumerate}

\begin{conjecture}[No dynamically coherent q4 shell descent]
\label{conj:shell}
No one same-trajectory full-defect q4 entrance with finite unweighted
dissipation can simultaneously have:
\begin{enumerate}
 \item positive spectral energy concentrated near
 \(L(s)=s^{-9/11}\ell(s)^{2/11}\);
 \item nonlinear transfer asymptotic to positive adjacent-shell Pell
 transfers at that moving shell; and
 \item a fixed fraction of every terminal energy quantum supplied from
 source times \(s=o(\tau)\).
\end{enumerate}
\end{conjecture}

\begin{remark}
Conjecture~\ref{conj:shell} is not proved.  The shell model is not a
velocity field and the Pell triads are snapshots.  No q4 closure,
regularity theorem, breakdown theorem, energy-equality theorem, or Clay
A--D resolution is asserted.
\end{remark}

\section*{Conclusion}

The extra positive-energy structure invalidates the previous Laguerre
countermodel but does not produce a contradiction.  A sharper moving
shell meets the complete-monotonicity and q4 requirements, and exact
quadratic triads reproduce its adjacent-shell heat shape.  The unknown
has therefore moved from radial algebra to dynamic compatibility:
whether one finite-energy NSE trajectory can repeatedly make those
transfers while its positive spectrum descends from infinity.  The Clay
problem remains unsolved.

\end{document}
